\begin{lemma}[Structural lemma; Scheide Lemma 4.1]
Let $G$ satisfy the source's critical-graph convention, let
$\chi'(G)=k+1$ with $k\geq\Delta(G)+1$, and let
$(T,e_0,\varphi)$ be a globally maximum triple carrying the exact ordered
S1--S4 block configuration of
\noderef{ScheideStructuralConfiguration}.  Then $G$ is elementary.
\end{lemma}
\begin{proof}
Apply \noderef{ScheideStructuralReduction} and replace the original ordered
tree by its reduced reordering $T'$.  This changes neither the coloring nor
the vertex set and retains every edge and vertex identity in the precise
way stated there.

If the reduced tree has no defective color, then
\noderef{ScheideGlobalMaximumBoundaryFacts} makes $V(T')$ strongly closed,
while \noderef{ScheideGlobalMaximumOrder} makes it elementary.  The first
edge of the tree-sequence is $e_0$, so its attachment equation puts both
ends of $e_0$ in $V(T')$.  The globally maximum field supplies criticality,
$\chi'(G)=k+1$, $k\geq\Delta(G)+1$, properness, and the fact that exactly
$e_0$ is uncolored; the configuration's no-isolated-vertex field supplies
the remaining source criticality convention.  Hence condition (iv) of
\noderef{egraph2} holds with $X=V(T')$, and its implication
(iv)$\Rightarrow$(i) yields that $G$ is elementary.

Assume now that a defective color exists.  Choose the incidence paths and
named departures supplied by \noderef{ScheideDefectiveSetupExists}, then
choose a replacement-data family by
\noderef{ScheideReplacementFamilyExists}.  Assertion (16),
\noderef{ScheideEnlargedElementary}, says that the resulting enlarged set
$X$ is elementary.  Assertion (21),
\noderef{ScheideEnlargedStronglyClosed}, says that the same set is strongly
closed.  It contains $V(T')$ by the definition
\noderef{ScheideEnlargedVertexSet}, and therefore contains both ends of
$e_0$.

The original proper coloring and unique-uncolored-edge clauses remain
unchanged, and the critical, chromatic-index, degree, and no-isolated-vertex
hypotheses are exactly those already listed in the no-defective case.
Thus condition (iv) of \noderef{egraph2} again holds, now with the enlarged
set $X$.  The implication (iv)$\Rightarrow$(i) proves that $G$ is
elementary.  This covers both cases and completes the structural lemma.
\end{proof}
